\documentclass[11pt,a4paper]{article}
\usepackage[utf8]{inputenc}
\usepackage{amsmath}
\usepackage{amssymb}
\usepackage{geometry}
\geometry{margin=1in}
\usepackage{hyperref}
\usepackage{graphicx}
\usepackage{booktabs}
\usepackage{caption}

\setcounter{secnumdepth}{0}

\hypersetup{
    colorlinks=true,
    linkcolor=blue,
    filecolor=magenta,      
    urlcolor=cyan,
}

\title{Ballooning Instability: Operating Principles of \texttt{Bal.jl} and the Perturbative Dynamics Behind the S-Alpha Diagram}
\author{}
\date{}

\begin{document}

\maketitle

This document discusses the physical and mathematical foundations of \texttt{Bal.jl},
which computes ballooning instability. It describes the full procedure from the
formulation of the ballooning-mode equation, through the matching of asymptotic
boundary conditions, to the perturbation framework used to generate the
\(s-\alpha\) diagram.

\begin{center}\rule{0.5\linewidth}{0.5pt}\end{center}

\section{1. Ballooning-Mode Equation and the Basic Operating Principle of \texttt{Bal.jl}}\label{ballooning-mode-equation-and-basic-operating-principle-of-baljl}

\subsection{1.1. Ballooning Representation}\label{ballooning-representation}

In a tokamak equilibrium, physical quantities must be periodic with respect to the
poloidal angle \(\theta\) and the toroidal angle \(\zeta\). Equation (1) is the
physical periodicity condition that the perturbation function \(\varphi\) must
satisfy. Here, \(q\) is the safety factor.

\[ \varphi(\psi, \theta, \beta) = \varphi(\psi, \theta, \beta + 1) = \varphi(\psi, \theta + 1, \beta - q) \tag{1} \]

Ballooning modes are characterized by very long wavelengths along magnetic field
lines and short wavelengths in the perpendicular direction. To describe this
behavior, the ballooning transformation in equation (2) is introduced using a
covering-space function \(\hat{\varphi}\), defined on the extended domain
\(-\infty < \theta < \infty\).

\[ \varphi(\psi, \theta, \beta) = \sum_{l=-\infty}^{\infty} \hat{\varphi}(\psi, \theta - l) \exp [-2\pi i n (\beta + lq)] \tag{2} \]

Here, \(\beta = \zeta - q\theta + \int q' \theta_0(\psi) d\psi\) is the Clebsch
coordinate. The magnitude of the wave-number vector \(\nabla\beta\) contains a
secular term proportional to \(\theta\), generated by the magnetic shear \(q'\), as
shown in equation (3).

\[ |\nabla\beta|^2 = |\nabla\zeta - q\nabla\theta|^2 + q'^2(\theta - \theta_0)^2 |\nabla\psi|^2 - 2q'(\theta - \theta_0)(\nabla\zeta - q\nabla\theta) \cdot \nabla\psi \tag{3} \]

\subsection{1.2. Ideal Marginal Ballooning Equation}\label{ideal-marginal-ballooning-equation}

At the stability boundary, \(\gamma=0\), variation of the energy principle yields
the ordinary differential equation in equation (4). Here, \(\chi\) is the poloidal
flux and \(P' = \mu_0 p\).

\[ \mathbf{B} \cdot \nabla \left( \frac{|\nabla\beta|^2}{B^2} \mathbf{B} \cdot \nabla \hat{\varphi} \right) + 2 \frac{P'}{\chi'} \kappa_w \hat{\varphi} = 0 \tag{4} \]

Written explicitly using the Jacobian \(\mathcal{J}\), this becomes equation (5).

\[ \frac{\chi'}{\mathcal{J}} \frac{\partial}{\partial\theta} \left( \frac{|\nabla\beta|^2}{B^2} \frac{\chi'}{\mathcal{J}} \frac{\partial \hat{\varphi}}{\partial\theta} \right) + 2 \frac{P'}{\chi'} \kappa_w \hat{\varphi} = 0 \tag{5} \]

To solve this numerically, \texttt{Bal.jl} defines the state vector \(\mathbf{u}\)
and converts the equation into the first-order matrix differential equation shown
in equation (6).

\[ \frac{\partial \mathbf{u}}{\partial\theta} = \mathbf{L}\mathbf{u}, \quad \mathbf{u} \equiv \begin{pmatrix} \hat{\varphi} \\ \frac{|\nabla\beta|^2}{B^2} \frac{\chi'}{\mathcal{J}} \frac{\partial \hat{\varphi}}{\partial\theta} \end{pmatrix}, \quad \mathbf{L} \equiv \frac{\mathcal{J}}{\chi'} \begin{pmatrix} 0 & \frac{B^2}{|\nabla\beta|^2} \\ -2\frac{P'}{\chi'} \kappa_w & 0 \end{pmatrix} \tag{6} \]

\subsection{1.3. Decomposition and Expansion of Curvature}\label{decomposition-and-expansion-of-curvature}

The key element that drives ballooning instability is the driving curvature. From
the MHD equilibrium condition \(\nabla P = \mathbf{J} \times \mathbf{B}\), the
magnetic-field-line curvature \(\boldsymbol{\kappa}\) is derived, and the combined
curvature \(\kappa_w\), projected along the wave-number vector direction, is
obtained. Substituting the Clebsch-coordinate relation
\(\nabla\beta = \nabla\zeta - q\nabla\theta - q'\theta\nabla\psi\) and expanding
it explicitly reveals the separation into a non-secular term \(\kappa_n\) and a
secular term \(-q'\theta\kappa_s\), as shown in equation (7).

\[ \begin{aligned} 
\kappa_w &\equiv \frac{\boldsymbol{\kappa} \times \mathbf{B}}{B^2} \cdot \nabla\beta = -\frac{1}{B} \nabla \cdot \left( \frac{\mathbf{B} \times \nabla\beta}{B} \right) \\ 
&= -\frac{1}{\mathcal{J} B} \left[ \frac{\partial}{\partial\psi} \left( \mathcal{J} \frac{\mathbf{B} \cdot \nabla\beta \times \nabla\psi}{B} \right) + \frac{\partial}{\partial\theta} \left( \mathcal{J} \frac{\mathbf{B} \cdot \nabla\beta \times \nabla\theta}{B} \right) \right] \\ 
&= -\frac{1}{\mathcal{J} B} \left[ \frac{\partial}{\partial\psi} \left( \mathcal{J} \frac{\mathbf{B} \cdot (\nabla\zeta - q\nabla\theta - q'\theta\nabla\psi) \times \nabla\psi}{B} \right) + \frac{\partial}{\partial\theta} \left( \mathcal{J} \frac{\mathbf{B} \cdot (\nabla\zeta - q\nabla\theta - q'\theta\nabla\psi) \times \nabla\theta}{B} \right) \right] \\ 
&= -\frac{1}{\mathcal{J} B} \left[ \frac{\partial}{\partial\psi} \left( \frac{\mathcal{J} B}{\chi'} \right) + \frac{\partial}{\partial\theta} \left( \mathcal{J} \frac{\mathbf{B} \cdot \nabla\zeta \times \nabla\theta}{B} - q'\theta \frac{\mathcal{J} (\mathbf{B} \cdot \nabla\psi \times \nabla\theta)}{B} \right) \right] \\ 
&= \left[-\frac{1}{\mathcal{J} B} \frac{\partial}{\partial\psi} \left( \frac{\mathcal{J} B}{\chi'} \right) + \frac{1}{\mathcal{J} B} \frac{\partial}{\partial\theta} \left( \mathcal{J} \frac{\mathbf{B} \cdot \nabla\theta \times \nabla\zeta}{B} \right) + \frac{2\pi f q'}{B^2 \mathcal{J}}\right] - \frac{q'\theta}{\mathcal{J} B} \frac{\partial}{\partial\theta} \left( \frac{\mathcal{J} \mathbf{B} \cdot \nabla\psi \times \nabla\theta}{B} \right) \\ 
&= \left[-\frac12 \left(\frac{1}{\chi'} \left( \frac{P'}{B^2} + \frac{\chi''}{\chi'} - \frac{\mathcal{J}'}{\mathcal{J}} \right) + \frac{2\pi f q'}{B^2 \mathcal{J}} + \frac{1}{\mathcal{J} } \frac{\partial}{\partial\theta} \left( \mathcal{J} \frac{\mathbf{B} \cdot \nabla\theta \times \nabla\zeta}{B^2} \right) \right)\right]- q'\theta\kappa_s \\ 
& =\kappa_n - q'\theta\kappa_s, \quad \left( \text{where } \kappa_s \equiv \frac{\boldsymbol{\kappa} \times \mathbf{B}}{B^2} \cdot \nabla\psi = -\frac{\pi f}{\mathcal{J} } \frac{\partial}{\partial\theta} \left( \frac{1}{B^2} \right) \right) 
\end{aligned} \tag{7} \]

This expansion and the final definition of the separated non-secular curvature
\(\kappa_n\) form the main backbone for assembling the differential coefficient
matrix.

\subsection{1.4. Large-\(\theta\) Limit and Canonical Transformation}\label{large-theta-limit-and-canonical-transformation}

In the limit \(\theta \to \infty\), the coefficient matrix diverges because
\(|\nabla\beta|^2 \simeq q'^2 \theta^2 |\nabla\psi|^2\). To handle this, the
canonical transformation matrix \(\mathbf{R} = \text{diag}(\theta^{-1/2}, \theta^{1/2})\)
is applied, yielding equation (8).

\[ \frac{\partial \mathbf{v}}{\partial\theta} = \mathbf{M}\mathbf{v}, \quad \mathbf{M} = \begin{pmatrix} \frac{1}{2\theta} & \frac{\mathcal{J}}{\chi'} \frac{B^2\theta}{|\nabla\beta|^2} \\ -\frac{\mathcal{J}}{\chi'} \frac{2P'\kappa_w}{\chi'\theta} & -\frac{1}{2\theta} \end{pmatrix} \tag{8} \]

To prevent divergence of the lowest-order matrix \(\mathbf{M}_{-1}\), the
\(\mathbf{S}\) matrix in equation (9) is applied, transforming the state space to
the periodic basis \(\mathbf{w}\).

\[ \mathbf{S} = \begin{pmatrix} 1 & 0 \\ \sigma & 1 \end{pmatrix}, \quad \sigma = -\frac{2\pi f P' q'}{\chi'^2 B^2} \tag{9} \]

For the new basis, the order matrices
\(\mathbf{N}_k \equiv \mathbf{S}^{-1}\mathbf{M}_k\mathbf{S}\) are used, and the
series recurrence relation in equation (10) follows.

\[ \frac{\partial \mathbf{w}^{(k+1)}}{\partial\theta} = [\mathbf{N}_0 - (\alpha+k)\mathbf{I}]\mathbf{w}^{(k)} + \sum_{l=1}^{k} \mathbf{N}_l \mathbf{w}^{(k-l)} \tag{10} \]

\subsection{1.5. Mercier Exponent and Big/Small Solutions}\label{mercier-exponent-and-bigsmall-solutions}

By the solvability condition, for \(k=0\),
\(\langle [\mathbf{N}_0 - \alpha\mathbf{I}] \mathbf{w}^{(0)} \rangle = 0\).
Therefore, the Mercier exponent \(\alpha\) is determined from the numerical
determinant equation in equation (11).

\[ \det[\langle \mathbf{N}_0 \rangle - \alpha \mathbf{I}] = 0 \implies \alpha = \pm \sqrt{-D_I} \tag{11} \]

The divergent basis, called the Big solution, and the decaying basis, called the
Small solution, are given by equation (12).

\[ \mathbf{u}_{large}^{(0)} \simeq \begin{pmatrix} \theta^{\alpha - 1/2} \\ \theta^{\alpha + 1/2} ( \sigma + w_{21} ) \end{pmatrix}, \quad \mathbf{u}_{small}^{(0)} \simeq \begin{pmatrix} \theta^{-\alpha - 1/2} \\ \theta^{-\alpha + 1/2} ( \sigma + w_{22} ) \end{pmatrix} \tag{12} \]

\subsection{1.6. Periodic Integral of the First-Order Correction and the Role of the First-Order Solvability Condition}\label{periodic-integral-of-first-order-correction-and-role-of-first-order-solvability-condition}

The zeroth-order solution in the simple ideal \(\theta \to \infty\) limit alone is
not sufficient to suppress numerical errors induced in the finite \(\theta_{max}\)
boundary region. To improve accuracy, \texttt{Bal.jl} uses the \(k=0\) state of
equation (10) and computes the derivative of the periodic component linearly.

\[ \frac{\partial \mathbf{w}^{(1)}}{\partial\theta} = (\mathbf{N}_0 - \alpha\mathbf{I})\mathbf{w}^{(0)} = \tilde{\mathbf{N}}_0(\theta)\mathbf{w}^{(0)} \tag{13} \]

The first-order functional solution obtained by integrating this relation consists
of a periodic function \(\tilde{\mathbf{w}}^{(1)}\) and an unknown constant matrix
\(\mathbf{c}^{(1)}\), as shown in equation (14).

\[ \mathbf{w}^{(1)}(\theta) = \left( \int \tilde{\mathbf{N}}_0(\theta) d\theta \right) \mathbf{w}^{(0)} + \mathbf{c}^{(1)} \equiv \tilde{\mathbf{w}}^{(1)}(\theta) + \mathbf{c}^{(1)} \tag{14} \]

To determine this constant structure vector, the solvability condition
corresponding to the \(k=1\) order is expanded, defining the closed inverse-matrix
equation for the integration constant estimate \(\mathbf{c}^{(1)}\), equation (15).

\[ [\langle \mathbf{N}_0 \rangle - (\alpha+1)\mathbf{I}] \mathbf{c}^{(1)} = - \left\langle [\mathbf{N}_0 - (\alpha+1)\mathbf{I}]\tilde{\mathbf{w}}^{(1)}(\theta) + \mathbf{N}_1 \mathbf{w}^{(0)} \right\rangle \tag{15} \]

The \texttt{spline\_integrate!} operation module in \texttt{Bal.jl} represents the
differential terms in equation (14) as splines, cumulatively integrates them over
one period, and mathematically computes the inverse matrix satisfying equation
(15) to obtain \(\mathbf{c}^{(1)}\) without gaps. Through this procedure, the
accurate first-order asymptotic boundary matching in equation (16), which includes
the first-order asymptotic term, is achieved.

\[ \mathbf{u}_{boundary} \simeq \mathbf{R}\mathbf{S} \left( \mathbf{w}^{(0)} + \frac{1}{\theta} \mathbf{w}^{(1)}(\theta) \right) \tag{16} \]

\begin{center}\rule{0.5\linewidth}{0.5pt}\end{center}

\section{2. Local Thin-Layer Ordering Framework}\label{local-thin-layer-ordering-framework}

The \(s-\alpha\) diagram is based on a theoretical framework that scans local
instability near a particular flux surface \(\psi_b\), without recomputing the full
background equilibrium. To formulate this rigorously, an asymptotic perturbation
parameter \(\mu\), with \(0 < \mu \ll 1\), is introduced and the thin-layer limit is
expanded.

\subsection{2.1. Stretched Coordinate}\label{stretched-coordinate}

To treat small variations around the target surface \(\psi_b\) on a macroscopic
scale, the radial coordinate is stretched into \(s\).

\[ s = \frac{\psi - \psi_b}{\mu} \tag{17} \]

Under this transformation, the radial partial-derivative operator is amplified to
order \(1/\mu\) through the \(s\) derivative.

\[ \frac{\partial}{\partial\psi} = \frac{1}{\mu} \frac{\partial}{\partial s} \tag{18} \]

\subsection{2.2. Local Perturbation Expansion of Equilibrium Quantities}\label{local-perturbation-expansion-of-equilibrium-quantities}

Basic profiles such as pressure \(P\), the flux derivative coefficient \(\chi'\),
and the safety factor \(q\) are expanded as the sum of a background equilibrium
(zeroth order) and a perturbed equilibrium (first order).

\[ P(\psi) = P^{(0)}(\psi) + \mu P^{(1)}(s) + \cdots \tag{19} \]
\[ q(\psi) = q^{(0)}(\psi) + \mu q^{(1)}(s) + \cdots \tag{20} \]

Here, the amplitude of each perturbation itself is considered to vary only within a
very small region of order \(O(\mu)\).

\subsection{2.3. First-Order Local Derivative Perturbations: Subscript \(s\)}\label{first-order-local-derivative-perturbations-subscript-s}

However, when these equilibrium quantities are differentiated with respect to
\(\psi\), for example \(P_\psi\), the chain rule cancels the amplitude factor
\(\mu\) against the reciprocal factor \(1/\mu\). As a result, strong variations of
order \(O(1)\) appear at the derivative-coefficient level.

\[ P_\psi = P_\psi^{(0)} + P_s^{(1)} \quad \left( \text{where } P_s^{(1)} \equiv \frac{\partial P^{(1)}}{\partial s} \right) \tag{21} \]
\[ q_\psi = q_\psi^{(0)} + q_s^{(1)} \quad \left( \text{where } q_s^{(1)} \equiv \frac{\partial q^{(1)}}{\partial s} \right) \tag{22} \]

Thus, terms carrying the subscript \(_s\) represent perturbative rates of change
promoted by the extreme-wavelength derivative to order \(O(1)\), comparable to the
zeroth-order background derivative terms \(P_\psi^{(0)}\) and \(q_\psi^{(0)}\).
These terms later become the practical variables, or axes, of the \(s-\alpha\)
diagram.

\subsection{2.4. Retained Contributions: Superscript \([1]\)}\label{retained-contributions-superscript-1}

When unit local perturbations such as \(P_s^{(1)}\) and \(q_s^{(1)}\) are combined
in the assembly of equations involving \(\mathbf{M}_0\), \(\mathbf{N}_0\), and other
matrices, the resulting first-order correction terms are collectively denoted by
the superscript \(^{[1]}\). A representative example is \(I_{per}^{[1]}\), the local
periodic variation of the integrated shear. This notation means that, in the
numerical limiting expansion \(\mu \to 0\), the term is retained all the way through
the equation system and represents a major physical effect, rather than being
truncated, linearized away, or removed.

\begin{center}\rule{0.5\linewidth}{0.5pt}\end{center}

\section{3. Introduction of \(P'\) and \(q'\) Perturbations for the \(s-\alpha\) Diagram}\label{introduction-of-pprime-and-qprime-perturbations-for-s-alpha-diagram}

Under this perturbation framework, the local perturbations \(P_s^{(1)}\),
\(q_s^{(1)}\), and \(I_{per}^{[1]}\) are combined without recomputing the
equilibrium. Using the integrated-shear function
\(I_{tot} = q'(\theta-\theta_k) + I_{per}\), the wave-number space and curvature are
expanded as equations (23) and (24), respectively.

\[ \left(\frac{|\nabla\beta|}{B}\right)^2 = \frac{1}{\chi'^2|\nabla\psi|^2} + \frac{|\nabla\psi|^2}{B^2}I_{\mathrm{tot}}^2 \tag{23} \]

\[ \kappa_w = \frac{\boldsymbol\kappa\cdot\nabla\psi}{\chi'|\nabla\psi|^2} - \kappa_s I_{\mathrm{tot}} \tag{24} \]

If the \(I_{tot}\) part of this expression is separated into the perturbation
expansion \(I_{tot} = I_{tot}^{(0)} + I_{tot}^{[1]}\), where
\(I_{tot}^{[1]} = q_s^{(1)}(\theta - \theta_k) + I_{per}^{[1]}\), the final retained
\(O(1)\) unit perturbation terms are obtained as equations (25) and (26).

\[ \left(\frac{|\nabla\beta|^2}{B^2}\right)_{O(1),\mathrm{ret}} = \frac{|\nabla\psi|_0^2}{B_0^2} \left[ 2I_{\mathrm{tot}}^{(0)}I_{\mathrm{tot}}^{[1]} + \bigl(I_{\mathrm{tot}}^{[1]}\bigr)^2 \right] \tag{25} \]

\[ \kappa_w^{[1]} = -\kappa_s^{(0)}I_{\mathrm{tot}}^{[1]} \tag{26} \]

The cross term and nonlinear term \(\bigl(I_{\mathrm{tot}}^{[1]}\bigr)^2\) derived
above are not discarded by a linearization process. Instead, they are retained in
the equation matrix with their original numerical order.

\begin{center}\rule{0.5\linewidth}{0.5pt}\end{center}

\section{4. Construction Framework of the \(\mathbf{N}_0\) Matrix}\label{construction-framework-of-the-n0-matrix}

The lower-triangular element \(\sigma\) of the canonical transformation matrix
\(\mathbf{S}\) is defined by combining the \(P'\) and \(q'\) perturbations, as in
equation (27).

\[ \sigma = -\frac{2\pi f}{\chi_0'^2 B_0^2} \bigl(P_\psi^{(0)}+P_s^{(1)}\bigr) \bigl(q_\psi^{(0)}+q_s^{(1)}\bigr) \tag{27} \]

The \((1,2)\) component of the asymptotic transformation matrix \(\mathbf{M}_0\)
responds only to variations in \(q_s^{(1)}\).

\[ (\mathbf{M}_0)_{12} = \frac{\mathcal{J}_0}{\chi_0'} \frac{B_0^2}{\bigl(q_\psi^{(0)}+q_s^{(1)}\bigr)^2|\nabla\psi|_0^2} \tag{28} \]

For the \((2,1)\) component of the coefficient matrix, the perturbation of the
non-secular curvature \(\kappa_n(\theta)\) is expanded as
\(\kappa_n(\theta) = \kappa_n^{(0)}(\theta) + \kappa_s^{(0)}(\theta)q_s^{(1)}\theta_k - \kappa_s^{(0)}(\theta)I_{per}^{[1]}(\theta)\),
with \(\theta_k\) treated rigorously. This leads to the form in equation (29).

\[ (\mathbf{M}_0)_{21} = -\frac{2\mathcal{J}_0}{\chi_0'^2} \bigl(P_\psi^{(0)}+P_s^{(1)}\bigr) \Bigl[ \kappa_n^{(0)}(\theta) + \kappa_s^{(0)}(\theta)q_s^{(1)}\theta_k - \kappa_s^{(0)}(\theta)I_{per}^{[1]}(\theta) \Bigr] \tag{29} \]

The final matrix
\(\mathbf{N}_0(\theta) = \mathbf{S}^{-1}\mathbf{M}_0(\theta)\mathbf{S}\) is then
integrated over one period to form \(\langle \mathbf{N}_0 \rangle\), which is used
in the solvability-condition eigenvalue equation (11).

\begin{center}\rule{0.5\linewidth}{0.5pt}\end{center}

\section{5. Correction of the \(\mathbf{N}_1\) Matrix for First-Order Asymptotic Numerical-Limit Matching}\label{correction-of-n1-matrix-for-first-order-asymptotic-numerical-limit-matching}

In the numerical code, to prevent divergence and achieve asymptotic finiteness over
a finite \(\theta_{max}\) range, the boundary condition includes not only the
zeroth-order solution but also the first-order asymptotic correction
\(\mathbf{w}^{(1)}(\theta)\), as in equation (16).

To compute this, the integration constants associated with the first-order
solvability conditions in equations (14) and (15) must be evaluated. This
inevitably requires the correction formula for \(\mathbf{N}_1\), the matrix
describing the \(O(\theta^{-2})\) behavior. In the large-\(\theta\) limit, extracting
the term proportional to \(\theta^{-2}\) from the wave-number-vector components
leaves only the upper-right component of \(\mathbf{M}_1\), giving equation (30).

\[ \mathbf{M}_1(\theta) = \begin{pmatrix} 0 & a_1(\theta) \\ 0 & 0 \end{pmatrix}, \quad a_1(\theta) = -2\frac{\mathcal{J}_0}{\chi_0'} \frac{B_0^2}{|\nabla\psi|_0^2\bigl(q_\psi^{(0)}+q_s^{(1)}\bigr)^3} \left[ I_{per}^{(0)}(\theta)+I_{per}^{[1]}(\theta) - \bigl(q_\psi^{(0)}+q_s^{(1)}\bigr)\theta_k \right] \tag{30} \]

Applying the \(\mathbf{S}^{-1}\) transformation once again gives the final
\(\mathbf{N}_1\) transmitted to the matching equation.

\[ \mathbf{N}_1(\theta) = a_1(\theta) \begin{pmatrix} \sigma & 1 \\ -\sigma^2 & -\sigma \end{pmatrix} \tag{31} \]

This formula also demonstrates that, apart from the background terms, the
\(\mathbf{N}_1(\theta)\) sequence is mathematically closed solely in terms of the
three local perturbation factors \(I_{per}^{[1]}(\theta)\), \(P_s^{(1)}\), and
\(q_s^{(1)}\). The terms \(I_{per}^{[1]}\) and \(q_s^{(1)}\) control the core
perturbation through \(a_1(\theta)\), while \(P_s^{(1)}\) governs the distribution
of coefficient space through the lower-triangular canonical element \(\sigma\).

In summary, this retained-order assignment control for the two fundamental
matrices \(\mathbf{N}_0\) and \(\mathbf{N}_1\) generates the boundary condition
\(\mathbf{u}_{boundary}\), whose basis solution does not collapse even near strict
equilibrium singularities such as edge effects. This constitutes the theoretical
foundation that enables \texttt{Bal.jl} to provide a highly reliable and internally
consistent local \(s-\alpha\) diagram scan.

\begin{center}\rule{0.5\linewidth}{0.5pt}\end{center}

\section{6. Practical Numerical Implementation in \texttt{Bal.jl}}\label{practical-numerical-implementation-in-baljl}

\subsection{6.1. Replacement of the Formula for Geometric Curvature Calculation}\label{replacement-of-formula-for-geometric-curvature-calculation}

The formula used to calculate the non-secular curvature \(\kappa_n\) is replaced by
equation (32), which uses only geometric information, rather than the previous
algebraic expansion that depended on the Grad--Shafranov (G--S) equation.

The G--S-dependent expression for \(\kappa_n\) in equation (7), used by the previous
code, assumes that the equilibrium satisfies the G--S equation. This expression has
computational advantages because it avoids some operations such as Jacobian partial
derivatives. However, equilibrium data obtained from actual numerical solvers
contain non-negligible G--S errors in several regions, including near the plasma
edge. As a result, unlike the mathematical expansion, mechanical numerical errors
accumulate and produce values that differ from the actual geometric magnetic-field
curvature.

Therefore, equation (32) is used, preserving the entanglement of geometric
information appearing in the derivation in Section 1.3 purely in spatial-derivative
form.

\[ \kappa_n^{(0)} = -\frac{1}{\mathcal{J}_0 B_0} \frac{\partial}{\partial\psi} \left( \frac{\mathcal{J}_0 B_0}{\chi'_0} \right) + \frac{1}{\mathcal{J}_0 B_0} \frac{\partial}{\partial\theta} \left( \mathcal{J}_0 \frac{\mathbf{B}_0 \cdot \nabla\theta \times \nabla\zeta}{B_0} \right) + \frac{2\pi f q'_0}{B_0^2 \mathcal{J}_0} \tag{32} \]

As a result, \texttt{Bal.jl} is not sensitively distorted by unavoidable G--S errors
in the input numerical equilibrium and achieves a much more robust,
geometry-dependent capability for calculating instability curvature.

\subsection{6.2. Dual-Loop Structure, Loop 1 and Loop 2, for Overcoming Geometric-Derivative Noise}\label{dual-loop-structure-for-overcoming-geometric-derivative-noise}

Summarizing the preceding expansion formulas, the full curvature expression can be
expanded so that it depends explicitly on \(P'\), \(q'\), \(I_{per}\), and related
quantities. Then the full \(s-\alpha\) scan can be performed easily by updating the
system with local perturbations such as \(P' = P'^{(0)} + P_s^{(1)}\).

However, a problem arises when numerically expanding and calculating the
wave-vector-coupled curvature \(\kappa_w\).

In Greene's original paper (1981), the curvature is intuitively decomposed as
\(\kappa_w = \kappa_n^{Greene} + \kappa_g\), where the normal curvature has the form
\(\kappa_n^{Greene} \propto \boldsymbol\kappa \cdot \nabla\psi\). For this definition
of \(\kappa_n^{Greene}\), the first-order perturbative contribution is effectively
zero, and the main perturbation dependence is concentrated in the geodesic
curvature \(\kappa_g\) multiplied by the integrated shear. If this decomposition
were followed directly, then subsequent perturbations would only need to be added
to \(\kappa_g\), which would be an ideal mathematical structure.

However, this decomposition differs from the secular/non-secular periodic
decomposition used in this document and in the existing code,
\(\kappa_w = \kappa_{n} - \kappa_s I_{tot}\). More importantly, if one attempts to
calculate \(\kappa_n^{Greene}\) using the intuitive inner-product expression
\(\boldsymbol\kappa \cdot \nabla\psi\) in Greene's approach, severe numerical errors
and noise blow up during differentiations of magnetic-field-structure quantities
inherent in numerical equilibrium data, such as \(\partial/\partial\psi(B^2)\). In
practice, this makes the method unusable within the system.

Therefore, to make the analysis robust, \texttt{Bal.jl} first computes the base
spatial-derivative expression for \(\kappa_n\), which requires numerical stability,
using the existing zeroth-order expansion method, equation (32), which is less
sensitive to geometric noise. Later, the first-order perturbative contribution to
the modified \(\kappa_w\) is separated algebraically and added separately. This
numerical bypass strategy is the fundamental reason the code was split into Loop 1
and Loop 2.

\begin{itemize}
\item \textbf{[Loop 1]}: Insert only the unperturbed, pure zeroth-order background
  parameters, such as \(P'\) and \(q'\), and safely precompute the base structural
  curvature \(\kappa_n^{(0)}\), equation (32), which contains geometric
  derivatives, without noise divergence.
\item \textbf{[Loop 2]}: Insert the local perturbed terms requested by the user,
  such as \(P_s^{(1)}\), \(q_s^{(1)}\), and \(I_{per}^{[1]}\), into the system. At
  this stage, unstable repeated differentiation of geometric parameters is omitted.
  Only the analytically derived algebraic first-order perturbation increment,
  \(\kappa_s^{(0)}q_s^{(1)}\theta_k - \kappa_s^{(0)}I_{per}^{[1]}\), is added to the
  previously computed \(\kappa_n^{(0)}\). The remaining terms, including
  \(\mathbf{M}_0\), are then computed sequentially and rapidly.
\end{itemize}

\subsection{6.3. Principle for Setting the Asymptotic Limit \(\theta_{max}\) and Adoption of the \texttt{Baloo.f} Matching Method}\label{principle-for-setting-theta-max-and-adoption-of-baloo-f-matching-method}

When obtaining the asymptotic solution of the ballooning mode or setting the
asymptotic limit in the large-\(\theta\) region, the previous ballooning-mode
analysis code, \texttt{bal.f}, used the formula
\(\theta_{max} = \sqrt{\max(\sqrt{c/a_2})} \times 10.0 \times \theta_{max0}\).

This formula originates from writing the quadratic-polynomial behavior of
\((\frac{\nabla \beta}{B})^2\) in \(\theta\) as a completed square.

\[ a_2\theta^2 + a_1\theta + a_0 = a_2 \left( \theta + \frac{a_1}{2a_2} \right)^2 + \left( a_0 - \frac{a_1^2}{4a_2} \right) \tag{34} \]

Here, the residual term is \(c = a_0 - a_1^2 / (4a_2)\). For the asymptotic
analysis to be mathematically valid, the \(\theta^2\) term must dominate the
constant residual \(c\). Numerically, this requires
\(\sqrt{|c/a_2|} \ll \theta\). Therefore, the system evaluates the scales
\(\text{ratio} = |c/a_2|\) over all \(\theta\) grids spanning the core and edge, and
uses their maximum value to define the integration boundary dynamically as
\(\theta_{max} = 10 \cdot \theta_{max0} \cdot \sqrt{\max(|c/a_2|)}\), the distance
at which asymptotic behavior is stably secured.

However, attempts to use the resulting \(\theta_{max}\) to compute the small and
big solutions and then extract the first-order coefficient \(c_{a1}\) or the
parameter \(\Delta'\) are numerically unstable.

As shown in Figure \ref{fig:delta_prime}, the implementation comparison data reveal
clear quality differences among methods when extracting \(\Delta'\) and related
coefficients:
\begin{figure}[h]
    \centering
    \includegraphics[width=0.8\textwidth]{sources/08.png}
    \caption{Comparison of numerical stability among methods for extracting \(\Delta'\) and \(c_{a1}\) (\texttt{psimax plot.png})}
    \label{fig:delta_prime}
\end{figure}

\begin{itemize}
\item \textbf{Computing \(c_{a1}\) using the Small solution}: the curve shape
  collapses completely and cannot be used for physical discrimination.
\item \textbf{Computing \(\Delta'\) using the Small solution}: the result has the
  qualitative form of a stability indicator, but it contains relatively large
  noise.
\item \textbf{\texttt{Baloo.f} method}: the complicated derivation of \(c_{a1}\) and
  the asymptotic Small-solution matching process are skipped entirely. Instead, the
  boundary condition is imposed by simply forcing the displacement to vanish,
  \(y(\pm\theta_{max}) \to 0\), at the limiting boundaries \(\pm\theta_{max}\), and
  \(\Delta'\) is computed.
\end{itemize}

The comparison shows that the \texttt{Baloo.f} method, which uses a simple boundary
replacement, is the simplest and numerically most stable approach, and it gives a
clear continuous distribution of \(\Delta'\), compared with the Small-solution
method that requires delicate matching in the asymptotic region. \texttt{Bal.jl}
adopts the \texttt{Baloo.f}-style calculation method.

\subsection{6.4. Ballooning Angle \(\theta_k\) and Integral-Gauge Correction}\label{ballooning-angle-theta-k-and-integral-gauge-correction}

The parameter \(\theta_k\) in the ballooning-mode expansion physically acts as an
integration constant specifying the location at which the total accumulated local
shear, namely the integrated magnetic shear \(I_{tot}\), is exactly zero. In other
words, as defined in Greene (1981) and related studies, the gauge condition
\(I_{tot}(\theta_k) = 0\) is adopted.

In \texttt{Bal.jl}, for computational convenience, \(I_{tot}\) is calculated by
separating it into a linearly increasing term \(q'\) and a periodic variation term
\(I_{per}\). Under the perturbation expansion, it is written as

\[ I_{tot} = \bigl(q'^{(0)} + q_s^{(1)}\bigr)(\theta - \theta_k) + \bigl(I_{per}^{(0)}(\theta) + I_{per}^{[1]}(\theta)\bigr) \tag{35} \]

To strictly satisfy Greene's \(I_{tot}(\theta_k) = 0\) gauge, substituting
\(\theta = \theta_k\) into the expression above must make the linear term vanish,
which requires \(I_{per}(\theta_k) = 0\).

The \(I_{per}\) functions generally have arbitrary initial reference values and
therefore do not pass through zero at \(\theta_k\). To correct this and align the
code exactly with the theory of the paper, the code evaluates the spline value at
\(\theta_k\) and subtracts that value over the entire interval as an offset
correction.

\[ I_{per}^{(0)}(\theta) \leftarrow I_{per}^{(0)}(\theta) - I_{per}^{(0)}(\theta_k) \]
\[ I_{per}^{[1]}(\theta) \leftarrow I_{per}^{[1]}(\theta) - I_{per}^{[1]}(\theta_k) \]

Through this gauge-correction procedure, the code operates in complete agreement
with the gauge system assumed by the theory, namely \(I_{tot}(\theta_k) = 0\), for
all equilibria and local perturbation state spaces.

\subsection{6.5. About \texttt{locstab\_fs}}\label{about-locstab-fs}

In conclusion, \texttt{Bal.jl} now computes both the Mercier criterion and the
ballooning-instability \(\Delta'\) internally. Accordingly, \texttt{Mercier.jl} can
be removed from the existing calculation routine. As shown in Figure
\ref{fig:mercier_comp}, the results from the previous \texttt{Mercier.jl} and the
new Mercier criterion calculation in \texttt{Bal.jl} agree perfectly. This
demonstrates that the physical and mathematical implementation of \texttt{Bal.jl}
has achieved the same level of reliability as the previously validated method.

\begin{figure}[h]
    \centering
    \includegraphics[width=0.8\textwidth]{sources/Mercier.png}
    \caption{Comparison of Mercier criterion calculations from \texttt{Mercier.jl} and \texttt{Bal.jl}}
    \label{fig:mercier_comp}
\end{figure}

In addition, although \(c_{a1}\) and \(c_{a2}\) were previously stored in
\texttt{locstab\_fs} entries 4 and 5, respectively, this has now been changed so
that \(\Delta'\) is stored in \texttt{locstab\_fs} entry 4.

The logic for drawing the \(s-\alpha\) diagram for a profile is written in
\texttt{Bal\_salpha\_delta\_di\_summary.ipynb}, which should be used as the
reference.

\end{document}
